\documentclass[11pt]{article}
\usepackage[margin=1in]{geometry}
\usepackage{amsmath,amssymb}
\usepackage{physics}
\usepackage{booktabs}
\usepackage{hyperref}

\newcommand{\eps}{\varepsilon}
\newcommand{\BZ}{\mathrm{BZ}}
\newcommand{\MBZ}{\mathrm{MBZ}}
\newcommand{\Nk}{N_k}
\newcommand{\Ang}{\text{\AA}}
\newcommand{\Imag}{\operatorname{Im}}
\newcommand{\kb}{\mathbf{k}}
\newcommand{\rb}{\mathbf{r}}
\newcommand{\Om}{\Omega}
\newcommand{\Bfield}{B}
\newcommand{\sgn}{\operatorname{sgn}}
\newcommand{\Phib}{\Phi_{\!\Om}}
\newcommand{\lB}{\ell_B}
\newcommand{\code}[1]{\texttt{#1}}

\title{Semiclassical Onsager Quantization with Berry-Phase Corrections\\
\large Conventions, signs, and numerical implementation}
\author{Notes for internal reference}
\date{}

\begin{document}
\maketitle

\begin{abstract}
\noindent
These notes fix every sign and normalization entering the semiclassical
Landau-level (LL) fan calculation, and document how the Berry curvature
that feeds it is computed.  The motivating problem is that the
Berry-phase term in the Onsager relation had, in our implementation,
both the wrong overall sign and the wrong parity in $\Bfield$.  The
parity error is invisible for one-sided field ranges and was found by
comparing a fan computed over $\delta\Bfield \in [-2,2]\,$T against the
exact Hofstadter spectrum; the sign error survived even that, and was
found only by scoring each field branch separately at four fluxes.  Both
are now corrected.  What remains unexplained is a flux-independent
factor of four in the right-hand side of the quantization condition,
carried as the parameter $\lambda$.
Sections~\ref{sec:bc}--\ref{sec:onsager}
are the physics; Sec.~\ref{sec:numerics} is implementation;
Sec.~\ref{sec:signs} is a summary table intended to survive into the
supplementary information.  Section~\ref{sec:open} lists what is
\emph{not} settled --- read it before quoting any of this in print.
\end{abstract}

\tableofcontents

%======================================================================
\section{Scope and notation}
%======================================================================

We consider a two-dimensional band structure $E_n(\kb)$ obtained either
at zero field (moir\'e plane-wave expansion) or as the Hofstadter
subbands of a system already at a background flux
$\Bfield_0$.  In the second case the semiclassical field variable is the
\emph{deviation} $\delta\Bfield$ from $\Bfield_0$, since $\Bfield_0$ is
already encoded in the subband dispersion and in its Berry curvature.
Throughout, $\Bfield$ denotes whichever field the semiclassical dynamics
actually sees; in the code this is the contents of \code{Blist}, and it
routinely takes both signs.

The carrier charge is $q = -e$ with $e>0$.  The flux quantum is
\begin{equation}
  \phi_0 = \frac{2\pi\hbar}{e},
  \qquad
  \lB = \sqrt{\frac{\hbar}{e\abs{\Bfield}}}.
\end{equation}
Band index $n$, LL index $\nu$.  $S(E)$ is the $\kb$-space area of a
closed constant-energy orbit.  $\Om_n(\kb)$ is the Berry curvature
($z$-component; the code calls it \code{Oz}) and $m_n(\kb)$ the orbital
magnetic moment (\code{Lz}, up to a prefactor collected below).

%======================================================================
\section{Berry curvature and orbital moment}\label{sec:bc}
%======================================================================

\subsection{Definitions}

With the Berry connection defined as
\begin{equation}\label{eq:connection}
  \mathbf{A}_n(\kb) = i\braket{u_n(\kb)}{\grad_\kb u_n(\kb)},
\end{equation}
the Berry curvature is $\Om_n = \curl_\kb \mathbf{A}_n$, whose
$z$-component is
\begin{equation}\label{eq:curvature_def}
  \Om_n^z
  = i\Bigl( \braket{\partial_{k_x} u_n}{\partial_{k_y} u_n}
          - \braket{\partial_{k_y} u_n}{\partial_{k_x} u_n} \Bigr)
  = -2\,\Imag \braket{\partial_{k_x} u_n}{\partial_{k_y} u_n}.
\end{equation}
The Berry phase accumulated on a closed circuit $C$ is
\begin{equation}\label{eq:berryphase}
  \gamma_C = \oint_C \mathbf{A}_n \cdot d\kb ,
\end{equation}
and by Stokes' theorem $\gamma_C = +\int_{S_C} \Om_n^z\, d^2k$
\emph{when $C$ is traversed counterclockwise}.  The orientation caveat is
the entire subject of Sec.~\ref{sec:parity}; note it now.

\subsection{Kubo (perturbative) form}

Inserting $\ket{\partial_\mu u_n} = \sum_{m\neq n}
\ket{u_m}\mel{u_m}{\partial_\mu H}{u_n}/(E_n - E_m)$ into
Eq.~\eqref{eq:curvature_def} and writing the velocity operator as
$v_\mu = \hbar^{-1}\partial_\mu H$ gives
\begin{equation}\label{eq:Oz_kubo}
  \boxed{\;
  \Om_n^z(\kb) = -2\hbar^2 \sum_{m \neq n}
  \frac{\Imag\bigl[ v^x_{nm}\, v^y_{mn} \bigr]}
       {(E_n - E_m)^2 + \eta^2}
  \;}
\end{equation}
and, for the orbital magnetic moment,
\begin{equation}\label{eq:Lz_kubo}
  \boxed{\;
  L_n^z(\kb) = \hbar^2 \sum_{m \neq n}
  \frac{(E_n - E_m)\,\Imag\bigl[ v^x_{nm}\, v^y_{mn} \bigr]}
       {(E_n - E_m)^2 + \eta^2}
  \;}
\end{equation}
with $v^\mu_{nm} = \mel{u_n}{v_\mu}{u_m}$ in the eigenbasis of
$H(\kb)$.  The physical orbital moment is
$m_n^z = (e/2\hbar)\,L_n^z$; the prefactor is applied where it is needed
(Sec.~\ref{sec:bfield}) rather than being folded into the stored array.

Three remarks:
\begin{enumerate}
\item Diagonal terms vanish identically. For Hermitian $v_x, v_y$ the
  products $v^x_{nn} v^y_{nn}$ are real, so $\Imag[\cdot]=0$ and the
  $m=n$ term can be included in the sum without harm --- which is what
  the code does.
\item $\eta$ (\code{eta\_kubo}, default $2\,$meV) regularizes the
  degenerate denominator.  It is a numerical convenience, not a physical
  broadening, and it is distinct from the moir\'e coupling $\eta$ and
  from the transport $\Gamma$.  Individual band Chern numbers are only
  meaningful when the band is isolated by gaps $\gg \eta$.
\item The sum over $m$ is truncated to a window of \code{nremotebands}
  states on either side of the target manifold.  Because the kernel
  decays as $1/(E_n-E_m)^2$ this converges quickly, but it is the same
  convergence issue as the \code{transport\_buffer} in the quantum
  transport calculation and should be checked, not assumed.
\end{enumerate}

\subsection{Velocity operators}

\subsubsection{Plane-wave (zero-field) basis}

For the moir\'e Hamiltonian expanded in plane waves, $H(\kb)$ depends on
$\kb$ analytically and
\begin{equation}
  v_\mu = \frac{1}{\hbar}\pdv{H}{k_\mu}.
\end{equation}
For the bilayer this is block-diagonal in the $\mathbf{Q}$-vector index:
each Dirac block contributes $v_F \sigma_\mu/\hbar$, and the trigonal
warping term $\gamma_3$ contributes an interlayer block
$v^x_{TB} = -v_3 U_2/\hbar$, $v^y_{TB} = +i v_3 U_2 /\hbar$.  The moir\'e
potential is $\kb$-independent in this basis and drops out of $v_\mu$.

\subsubsection{Landau-level (Hofstadter) basis}\label{sec:vel_LL}

In the LL basis the situation is subtler and is worth spelling out,
because the naive route fails.  The physical velocity is
\begin{equation}
  v = \frac{i}{\hbar}\,\comm{R}{H},
\end{equation}
with $R$ the position operator.  In the magnetic Bloch representation
$R = i\,\partial_\kb + \mathbf{A}$, where $\mathbf{A}$ is the LL Berry
connection, so
\begin{equation}\label{eq:vel_general}
  v = \frac{1}{\hbar}\Bigl( \pdv{H}{\kb} + i \comm{\mathbf{A}}{H} \Bigr).
\end{equation}
Now split $H = H_{\text{base}} + V_{\text{moir\'e}}$.  The moir\'e
potential is a local function of position, so $\comm{R}{V_{\text{moir\'e}}}
= 0$, i.e.
\begin{equation}\label{eq:Vmoire_cancel}
  \pdv{V_{\text{moir\'e}}}{\kb} + i\comm{\mathbf{A}}{V_{\text{moir\'e}}} = 0
  \qquad\text{identically.}
\end{equation}
The BLG kinetic Hamiltonian $H_{\text{base}}$ carries no $\kb$
dependence, so $\partial H_{\text{base}}/\partial \kb = 0$ and
Eq.~\eqref{eq:vel_general} collapses to
\begin{equation}\label{eq:vel_LL}
  \boxed{\; v = \frac{i}{\hbar}\comm{\mathbf{A}}{H_{\text{base}}} \;}
\end{equation}
which is $\kb$-independent and can be precomputed once.

\paragraph{Why not evaluate the two terms separately.}
Equation~\eqref{eq:Vmoire_cancel} is an operator identity that requires
\emph{completeness} of the LL basis.  In a basis truncated at $N$ it
fails near the cutoff: $\partial V/\partial\kb$ and
$i\comm{\mathbf{A}}{V}$ each have large matrix elements connecting to
$\mathrm{LL}_N$, and their cancellation is spoiled.  Computing them
separately and subtracting therefore injects truncation artifacts into
$\Om_z$ precisely where the basis is least trustworthy.
Equation~\eqref{eq:vel_LL} sidesteps this by never forming either term.

\paragraph{The LL Berry connection.}
$\mathbf{A}$ connects adjacent LLs through the ladder operators.  Within
each sublattice ($s = A,B$) and each guiding-centre index, the nonzero
upper-triangular elements are
\begin{equation}\label{eq:ladder}
  \mel{s,\nu}{A_x}{s,\nu+1} = -\,i\,\frac{\lB}{\sqrt{2}}\sqrt{\nu+1},
  \qquad
  \mel{s,\nu}{A_y}{s,\nu+1} = +\,\frac{\lB}{\sqrt{2}}\sqrt{\nu+1},
\end{equation}
and the matrices are then Hermitized, $\mathbf{A} \to \mathbf{A} +
\mathbf{A}^\dagger$.  Both valleys use the same expression; they differ
only in which sublattice's top LL is chopped (Sec.~\ref{sec:chop}).

\paragraph{Gauge invariance.}
$\Om_z$ built from Eq.~\eqref{eq:vel_LL} is independent of the choice of
magnetic unit cell (square vs.\ triangular construction), which is a
useful check: the Chern number per mesh cell comes out as $1/pp$ and
integrates to an integer over the magnetic BZ.

\subsubsection{Spurious-mode chopping}\label{sec:chop}

In each valley one sublattice's highest LL has no coupling partner in
the off-diagonal kinetic term, producing a spurious zero mode.  It is
removed by deleting the corresponding rows and columns: the $(B,
\mathrm{LL}_N)$ states in valley $K$, the $(A, \mathrm{LL}_N)$ states in
valley $K'$.  The same deletion is applied to $H$, to $\mathbf{A}$, and
inside the interlayer assembly, so all operators remain
dimension-consistent.

\subsection{Units}

Internally $\Om_z$ is accumulated in $\Ang^2$ and $L_z$ in
$\mathrm{eV}\,\Ang^2$; on output they are converted with
$\Om_z \times 10^{-20} \to \mathrm{m}^2$ and
$L_z \times 10^{-20}\times 10^{3} \to \mathrm{meV}\,\mathrm{m}^2$.
Energies are parsed in meV, carried internally in eV, and returned in
meV.  The $\eta$ in Eqs.~\eqref{eq:Oz_kubo}--\eqref{eq:Lz_kubo} is
converted from meV to eV at read time.

%======================================================================
\section{The Onsager relation}\label{sec:onsager}
%======================================================================

\subsection{Semiclassical dynamics}

The wavepacket equations of motion in band $n$ are
\begin{align}
  \dot{\rb} &= \frac{1}{\hbar}\pdv{E_n}{\kb} - \dot{\kb}\times \Om_n, \\
  \hbar\dot{\kb} &= -e\,\dot{\rb}\times \mathbf{B},
  \qquad \mathbf{B} = \Bfield\hat{z}.
  \label{eq:eom_k}
\end{align}
Writing Eq.~\eqref{eq:eom_k} in components,
\begin{equation}\label{eq:eom_components}
  \hbar \dot{k}_x = -e\Bfield\,\dot{y},
  \qquad
  \hbar \dot{k}_y = +e\Bfield\,\dot{x}.
\end{equation}
The orbit is a constant-energy contour; what matters below is the
\emph{sense} in which it is traversed, which Eq.~\eqref{eq:eom_components}
shows is reversed under $\Bfield \to -\Bfield$.

\subsection{Quantization condition}

Onsager--Lifshitz quantization of the cyclotron orbit, including the
Berry phase $\gamma_C$ of Eq.~\eqref{eq:berryphase} and the Maslov index
$1/2$, reads
\begin{equation}\label{eq:onsager_textbook}
  S(E)\,\lB^2 = 2\pi\Bigl( \nu + \tfrac12 - \frac{\gamma_C}{2\pi} \Bigr),
  \qquad \nu = 0,1,2,\dots
\end{equation}
or equivalently, using $\lB^{-2} = e\abs{\Bfield}/\hbar$ and
$\phi_0 = 2\pi\hbar/e$,
\begin{equation}\label{eq:onsager_phi0}
  \frac{S(E)}{(2\pi)^2}
  = \frac{\abs{\Bfield}}{\phi_0}
    \Bigl( \nu + \tfrac12 - \frac{\gamma_C}{2\pi} \Bigr).
\end{equation}
Equation~\eqref{eq:onsager_phi0} is the form the code solves, and every
correction term below is a additive contribution inside the bracket.

\subsection{The Berry-phase term and its parity in \texorpdfstring{$\Bfield$}{B}}
\label{sec:parity}

This is the delicate part.  Define the \emph{enclosed Berry flux}
\begin{equation}\label{eq:Phib}
  \Phib(E) \;=\; \int_{S_C} \Om_n^z(\kb)\, d^2k ,
\end{equation}
the integral of the curvature over the orbit \emph{interior}.  Note
carefully what $\Phib$ is and is not:
\begin{itemize}
\item It is a property of the region enclosed by the contour.  It has
  \emph{no orientation attached}: the interior is a set of points, not a
  circuit.
\item It is therefore \emph{not} the Berry phase.  The Berry phase
  $\gamma_C$ is a line integral along a directed circuit, and by Stokes'
  theorem $\gamma_C = \pm\Phib$ with the sign set by the traversal sense.
\end{itemize}
Since Eq.~\eqref{eq:eom_components} reverses the traversal sense under
$\Bfield \to -\Bfield$ while Eq.~\eqref{eq:Phib} is manifestly
independent of $\Bfield$'s sign at fixed orbit, we have the central
structural result:
\begin{equation}\label{eq:parity}
  \boxed{\;\gamma_C = +\,\sgn(\Bfield)\,\Phib
  \quad\Longrightarrow\quad
  \gamma_C \text{ is \emph{odd} in } \Bfield. \;}
\end{equation}
Every other term in the bracket of Eq.~\eqref{eq:onsager_phi0} is even in
$\Bfield$.  The \emph{parity} is not a convention but a consequence of
Eq.~\eqref{eq:eom_components}, and it is what the numerical evidence most
sharply confirms.  The overall sign follows from the traversal sense: for
an electron-like orbit at $\Bfield > 0$, $\dot{\rb} \propto \grad_\kb E$
points outward and Eq.~\eqref{eq:eom_components} gives $\dot\kb \propto
(-k_y, k_x)$, i.e.\ counterclockwise, so $\gamma_C = +\Phib$ by Stokes.
It is independently confirmed by the exact Hofstadter spectrum at four
fluxes (Sec.~\ref{sec:validation}).  Earlier versions of this note and of
the code carried a minus here, in conflict with the traversal argument;
that conflict is resolved in favour of Eq.~\eqref{eq:parity} as written.

Two corollaries that are easy to conflate and must not be:
\begin{description}
\item[The sign is not per-orbit.] It follows from the carrier charge and
  the field direction, both global.  It does \emph{not} flip between
  electron-like and hole-like orbits, and it must never be inferred from
  the winding of a numerically extracted contour.
\item[The sign is per-field-direction.] A field list straddling zero
  requires both branches within a single calculation.
\end{description}
$\gamma_C$ is discontinuous at $\Bfield = 0$.  This is harmless: the
right-hand side of Eq.~\eqref{eq:onsager_phi0} vanishes there and no
Landau level is defined.

\subsection{Orbital moment and susceptibility corrections}

Two further corrections are carried, both even in $\Bfield$ inside the
bracket as currently implemented:
\begin{equation}
  \text{(orbital moment)}\quad -\frac{1}{2\pi}\dv{L}{E},
  \qquad
  \text{(Fukuyama susceptibility)}\quad
  -\,2\pi\,\dv{\chi}{E}\,\frac{\Bfield}{\phi_0}.
\end{equation}
The $\chi$ term descends from a $\Bfield^2$ energy correction, so its
linear-in-$\Bfield$ appearance in the bracket is expected.  The
$dL/dE$ term is flagged in Sec.~\ref{sec:open}.

\subsection{Form as implemented}

Collecting everything, with dimensionless prefactors $c_\Om, c_L,
c_\chi$ exposed as the input parameter \code{term\_factors}:
\begin{equation}\label{eq:implemented}
  \boxed{\;
  \frac{S(E)}{(2\pi)^2}
  = \frac{\abs{\Bfield}}{\phi_0}
  \left[
    \lambda\Bigl(\nu + \tfrac12\Bigr)
    - c_\Om\,\frac{\sgn(\Bfield)\,\Phib}{2\pi}
    - c_L\,\frac{1}{2\pi}\dv{L}{E}
    - c_\chi\, 2\pi \dv{\chi}{E}\frac{\Bfield}{\phi_0}
  \right]
  \;}
\end{equation}
Comparing with Eq.~\eqref{eq:onsager_phi0} identifies
$\gamma_C = +\sgn(\Bfield)\,c_\Om \Phib$, consistent with
Eq.~\eqref{eq:parity} at the default $c_\Om = 1$.

In the source, the residual is assembled as
$\mathrm{rhs} + \mathrm{base}$ with
$\mathrm{rhs} = \lambda\Bfield(\nu+\tfrac12)/\phi_0$ and \code{base\_S}
carrying an explicit $-\sgn(\Bfield)$; dividing through by
$\sgn(\Bfield)$ is what puts every term on $\abs{\Bfield}$ and produces
Eq.~\eqref{eq:implemented}.  The $\sgn(\Bfield)$ on the curvature term is
implemented as $\abs{\Bfield}$ in \code{base\_SB} where all other terms
use $\Bfield$ --- a one-character distinction that carries the whole of
Eq.~\eqref{eq:parity}'s parity.  Its overall sign is the leading minus on
that same line.

The multiplier $\lambda$ in Eq.~\eqref{eq:implemented} is the input
parameter \code{onsager\_Bmultiplier}, applied as
$\mathrm{rhs} \to \lambda\,\Bfield(\nu+\tfrac12)/\phi_0$.  It defaults to
$\lambda = 4$, for the reasons given in Sec.~\ref{sec:lambda}.  It does
\emph{not} enter the modified dispersion of Sec.~\ref{sec:bfield}, so it
cannot affect the band structure or $\Om^z$.

\subsection{The multiplier \texorpdfstring{$\lambda$}{lambda}}
\label{sec:lambda}

With $\lambda = 1$ --- the form Eq.~\eqref{eq:onsager_phi0} actually
derives --- the computed fan is too sparse by roughly a factor of four
relative to the exact Hofstadter spectrum.  Scanning $\lambda$ against
that spectrum puts the optimum at $4$ at every flux and geometry tested:

\begin{center}
\begin{tabular}{lccccc}
\toprule
$q/p$ & $p,q$ & unfolded & $v_M/A_{uc}$ & $2p$ & $\lambda_{\rm opt}$ \\
\midrule
$1/2$ & 2,1 & no  & 4.00 & 4  & 4 \\
$1/3$ & 3,1 & no  & 9.00 & 6  & 3--4 \\
$2/5$ & 5,2 & yes & 6.25 & 10 & 4 \\
$2/3$ & 3,2 & yes & 2.25 & 6  & 4 \\
\bottomrule
\end{tabular}
\end{center}

\noindent
$\lambda \approx 4$ holds while $v_M/A_{uc}$ ranges over $2.25$--$9$ and
$2p$ over $4$--$10$, folded and unfolded.  Both of those geometric
candidates were tested explicitly and score worse than $4$ at $2/5$ and
$1/3$.  A flux-independent constant is not a geometry factor: it is a
missing numerical factor in the right-hand side of
Eq.~\eqref{eq:implemented}, or equivalently in the normalization of
$S(E)$.  Two independent indicators single out $\lambda = 4$: the level
count comes out at about one semiclassical level per exact subband
($n_{\rm LL}/n_{\rm ex} \approx 1.05$ at $2/3$), and the sign test of
Sec.~\ref{sec:validation} discriminates by a factor of $\sim\!4$, whereas
at $\lambda = 2,3,5,6$ the two candidate signs score within ${\sim}30\%$
of each other with no consistent pattern.

Where the factor of four originates has not been traced; see
Sec.~\ref{sec:open}.  Until it is, it is carried explicitly on the
right-hand side rather than absorbed into $S(E)$, so that the place it is
being applied remains visible.

%======================================================================
\section{Numerical implementation}\label{sec:numerics}
%======================================================================

\subsection{Mesh and normalization}

The magnetic BZ is sampled on an $n_{k1}\times n_{k2}$ mesh with
$\Nk = n_{k1}n_{k2}$.  Flattening uses Fortran (column-major) order
throughout, matching the original MATLAB implementation.  With
$\mathrm{vol}_M$ the real-space moir\'e cell area,
\begin{equation}
  A_\BZ = \frac{(2\pi)^2}{\mathrm{vol}_M},
  \qquad
  a_{\text{cell}} = \frac{A_\BZ}{\Nk},
  \qquad
  w_k \equiv \frac{(2\pi)^2}{\Nk\,\mathrm{vol}_M} = a_{\text{cell}} .
\end{equation}
In Hofstadter mode $\mathrm{vol}_M = pp^2\,A_{uc}/(2qq)$ with
$A_{uc} = \sqrt{3}L^2 = 2A_{\text{prim}}$ the doubled (rectangular)
construction cell; equivalently
$A_\BZ = qq\,(2\pi)^2/(pp^2 A_{\text{prim}})$.  The flux per
\emph{primitive} moir\'e cell is $qq/(2pp)$ flux quanta, \emph{not}
$qq/pp$.

\subsection{Orbit detection}

For each band and each energy level $E_i$:
\begin{enumerate}
\item Reshape $E(\kb)$ to the 2D mesh and tile it $3\times 3$ so that
  orbits crossing the zone boundary close properly.
\item Run marching squares (\code{skimage.measure.find\_contours}) at
  level $E_i$.
\item Discard open contours (endpoints separated by more than one pixel)
  and contours with fewer than 4 vertices; close the remainder.
\item Deduplicate the $3\times3$ tiling by keeping only contours whose
  centroid lies in the central tile.
\item Compute the area by the shoelace formula and convert
  pixel$^2 \to k$-space via $a_{\text{cell}}$.
\item Reject areas below $a_{\text{cell}}$ (mesh noise) or above
  $A_\BZ - a_{\text{cell}}$ (a contour that has wrapped the whole zone).
\item Identify interior mesh points with
  \code{matplotlib.path.Path.contains\_points}, map tiled indices back
  to the primitive zone by modulo, and linearize in F-order.
\end{enumerate}
Multiple pockets at the same energy are retained and sorted by
descending area; the solver uses the dominant (largest) pocket.

\subsection{Enclosed Berry flux}

Equation~\eqref{eq:Phib} is evaluated as a plain mesh sum over the
interior points identified above,
\begin{equation}\label{eq:Phib_numeric}
  \Phib(E_i) = w_k \sum_{\kb \,\in\, \text{interior}} \Om_n^z(\kb),
\end{equation}
and the orbital-moment derivative as a Fermi-broadened mesh sum,
\begin{equation}
  \dv{L}{E}(E_i) = w_k \sum_{\kb} f'(E_n(\kb) - E_i)\, L_n^z(\kb),
  \qquad
  f'(\Delta) = \frac{e^{-\abs{x}}}{k_BT\,(1+e^{-\abs{x}})^2},
  \quad x = \frac{\Delta}{k_BT},
\end{equation}
with $k_BT$ the input parameter \code{kT} (default $3\,$meV).  The
$\abs{x}$ form is the numerically stable way to write
$\mathrm{sech}^2$.

\subsection{Orientation independence, and why it matters}
\label{sec:orientation}

Both of the quantities extracted above are orientation-blind by
construction:
\begin{itemize}
\item \code{\_shoelace\_area} takes $\abs{\cdot}$ of the signed shoelace
  sum, so the area is positive regardless of winding.
\item The interior is selected by point-in-polygon containment, not by a
  winding number, so Eq.~\eqref{eq:Phib_numeric} is likewise
  orientation-blind.
\end{itemize}
The orientation returned by marching squares is in fact mixed --- a
census over a single band finds both senses at different energies, with
no physical meaning: it reflects the contouring algorithm's internal
bookkeeping, not the direction of semiclassical motion.

The consequence is important and is the design decision worth stating
explicitly in any writeup: \textbf{the traversal sense is supplied
analytically, as the $\sgn(\Bfield)$ of Eq.~\eqref{eq:parity}, and is
never read off the numerics.}  Attempting to recover it from contour
winding would be both wrong (the winding is an artifact) and
unnecessary (the physical answer is a global function of $\sgn\Bfield$).

\subsection{Lifshitz segmentation}

When $S(E)$ is non-monotonic --- small pockets growing through a
Lifshitz transition into a large orbit that then shrinks --- the
quantization condition can have several roots at the same
$(\Bfield,\nu)$.  Since the solver returns one root per $(\Bfield,\nu)$,
the energy grid is first split into monotonic segments.

Detection is by magnitude: a transition is flagged wherever
$\abs{\Delta S}$ between adjacent grid points exceeds
$\texttt{lifshitz\_threshold}$ (default $50$) times the median
$\abs{\Delta S}$.  Each segment is solved independently and output keys
acquire a \code{\_seg}$i$ suffix.  This heuristic works but is
admittedly arbitrary; splitting instead at sign changes of $dS/dE$ would
be parameter-free and more physical, at the cost of sensitivity to noise
in $S(E)$.

\subsection{Root finding}

For each $(\Bfield, \nu)$ the residual of Eq.~\eqref{eq:implemented} is
evaluated on the segment's energy grid and rooted by:
\begin{enumerate}
\item \textbf{Sign-change interpolation} (primary): locate the first
  adjacent pair of grid points bracketing a sign change and interpolate
  linearly.  This gives continuous $E_\nu(\Bfield)$ curves even where
  $S(E)$ varies steeply, because the root need only fall \emph{between}
  grid points.
\item \textbf{Thresholded \code{argmin}} (fallback): where no sign change
  exists, take the grid point of smallest $\abs{\text{residual}}$ and
  reject it as \code{NaN} if that residual exceeds
  $r_{\text{tol}} = 5\%$ of $\abs{\Bfield}(\nu+\tfrac12)/\phi_0$.  This
  suppresses spurious levels at saddle energies where $S(E)$ has a
  maximum but the condition is never actually satisfied.
\end{enumerate}
Energies with zero orbit area are assigned an infinite residual so they
can never be selected --- a deliberate difference from the original
MATLAB implementation, which evaluated the condition at all energies and
could pick up spurious minima driven by the $dL/dE$ term in regions with
no orbit at all.

Grid resolution matters: with $k_BT = 3\,$meV broadening, energy grids
coarser than ${\sim}2\,$meV produce one-bin LL shifts.  A dedicated
denser grid (\code{elist\_onsager}) is available for this reason.

%======================================================================
\section{Non-perturbative field channel}\label{sec:bfield}
%======================================================================

The perturbative treatment folds the orbital moment in through
$dL/dE$.  An alternative --- and the channel used for the Hofstadter
comparisons --- is to put it directly into the dispersion.  At each
field,
\begin{equation}\label{eq:Emod}
  E^{\text{mod}}_n(\kb;\Bfield)
  = E_n(\kb) + g\,\Bfield\,\frac{e}{2\hbar}\,L_n^z(\kb),
\end{equation}
with $g$ the input parameter \code{gfactor}.  Contours, areas and
enclosed fluxes are then recomputed on $E^{\text{mod}}$ at every field
value, and the quantization condition is solved with $dL/dE \equiv 0$.

Consequences worth noting:
\begin{itemize}
\item Output labels change: the two cumulative sets are \code{\_SM}
  (area only, moment already in the dispersion) and \code{\_SBM} (plus
  enclosed curvature).  The solver's own \code{S}/\code{SB} outputs are
  renamed accordingly.
\item \code{term\_factors} is read as \emph{two} elements here,
  $[c_\Om\; c_\chi]$, and expanded internally to $(c_\Om, 0, c_\chi)$.
  A third element is silently ignored, since $c_L$ would multiply an
  identically zero array.
\item Equation~\eqref{eq:Emod} is manifestly odd in $\Bfield$, as it
  must be, and \code{onsager\_Bmultiplier} does \emph{not} appear in it.
\item Because contours must be re-found at every field, this channel is
  far more expensive than the perturbative one.  Intermediate data ---
  \code{Blist}, and per-band \code{area}, \code{enclosedBC},
  \code{E\_levels} arrays --- are therefore written to a separate
  \code{*\_detail.mat}.  These four arrays are the complete input to the
  quantization step, so the fan can be re-solved from them in seconds
  with different prefactors.  This is how the sign question was settled.
\end{itemize}

%======================================================================
\section{Validation}\label{sec:validation}
%======================================================================

Three checks, addressing different pieces.

\paragraph{(a) $\Om_z$ itself, via $\sigma_{xy}$.}
Semiclassical Chern numbers $C_n = w_k\sum_\kb \Om_n^z/2\pi$ were
compared against the plateau steps of the exact Kubo $\sigma_{xy}(\mu)$
from the quantum Hofstadter calculation at $qq/pp = 1/2$.  Bands 5--11
map one-to-one onto quantum subbands with $C_n = \Delta\sigma_{xy}$
agreeing to three decimals in every case.  This fixes the sign and
normalization of Eq.~\eqref{eq:Oz_kubo}.  Note that it constrains
$\Om_z$, \emph{not} the sign of the Berry term in the quantization
condition, so it is silent on Eq.~\eqref{eq:parity}.

\paragraph{(b) Overall sign, via a fixed-parity scan on both branches.}
This is the determination.  The recomputation utility can switch the
$\Bfield$-parity of the curvature term off, applying the same $c_\Om$ on
both field branches; because $\Bfield = \sgn(\Bfield)\abs{\Bfield}$, this
is achieved by handing the solver $c_\Om \sgn(\Bfield)$ per field point,
with no change to the solver itself.  Each branch can then be scored
against the exact Hofstadter spectrum independently.  At all four fluxes
tested --- $qq/pp = 1/2,\,1/3,\,2/5,\,2/3$, folded and unfolded --- the
spectrum selects
\begin{equation*}
  c_\Om = +1 \ \text{on}\ \delta\Bfield < 0,
  \qquad
  c_\Om = -1 \ \text{on}\ \delta\Bfield > 0,
\end{equation*}
which is the single odd-in-$\Bfield$ convention of
Eq.~\eqref{eq:parity} with $c_\Om = 1$.  The discriminating feature is
not the residual magnitude but the \emph{crossover}: the two curves swap
ranking exactly at the background field, which a fitting artifact would
not do.  Both spike identically at the background field itself, as
expected --- the right-hand side is proportional to $\delta\Bfield$, so
all $\nu$ collapse to one energy there.  Margins are $3.6\times$ and
$3.9\times$ at $1/2$ and $3.8\times$ and $4.9\times$ at $2/3$, the latter
with median residuals of $0.099$ and $0.083$\,meV.  The scan is only
meaningful at $\lambda = 4$ (Sec.~\ref{sec:lambda}).

\paragraph{(c) Parity, self-consistency.}
At $qq/pp = 1/3$ over $\delta\Bfield \in [-2,2]\,$T (200 fields, 13
bands, both valleys), the fan solved with Eq.~\eqref{eq:implemented}
reproduces the $c_\Om = +1$ result for \emph{every} $\delta\Bfield > 0$
and the $c_\Om = -1$ result for \emph{every} $\delta\Bfield < 0$: 1436 LL
arrays, 25\,386 finite entries, $\max\abs{\text{diff}} = 0$, no
\code{NaN}-pattern mismatches.  This confirms that the parity is plumbed
through the code as intended; it is not independent physical evidence.
The same splice was verified on a synthetic band with analytically
prescribed $S(E)$ and $\Phib(E)$.

\paragraph{Historical note.}
Two things made the earlier, incorrect sign hard to see.  A one-sided
field range hides the parity entirely: a plot window of
$[6.1, 11.8]\,$T against a background $\Bfield_0 = 11.555\,$T
corresponds to $\delta\Bfield \in [-5.45, +0.25]\,$T --- essentially
all negative --- so a uniform $c_\Om = -1$ appeared to work everywhere.
And a qualitative overlay, even a two-sided one, does not separate the
candidate signs; only the per-branch quantitative scoring of check (b)
does.  Fan files produced before the correction are stale in \code{\_SB},
\code{\_SBM} and \code{\_SBMC}.

%======================================================================
\section{Summary of sign conventions}\label{sec:signs}
%======================================================================

\begin{table}[h]
\centering
\begin{tabular}{@{}lll@{}}
\toprule
Quantity & Convention & Where fixed \\
\midrule
Berry connection & $\mathbf{A} = i\braket{u}{\grad_\kb u}$
  & Eq.~\eqref{eq:connection} \\
Berry curvature & $\Om^z = -2\hbar^2\sum_m
  \Imag[v^x_{nm}v^y_{mn}]/(\Delta E^2+\eta^2)$
  & Eq.~\eqref{eq:Oz_kubo}; validated by $\sigma_{xy}$ \\
Chern number & $C = w_k\sum_\kb \Om^z/2\pi$, $\sigma_{xy}=C\,e^2/h$
  & check (a) \\
Enclosed flux & $\Phib = \int_{\text{interior}}\Om^z d^2k$, no orientation
  & Eq.~\eqref{eq:Phib}, Sec.~\ref{sec:orientation} \\
Orbit area & $S = \abs{\text{shoelace}}$, always positive
  & Sec.~\ref{sec:orientation} \\
Berry phase & $\gamma_C = +\sgn(\Bfield)\,\Phib$ \quad(\emph{odd in} $\Bfield$)
  & Eq.~\eqref{eq:parity}; traversal $+$ check (b) \\
Maslov index & $+\tfrac12$ & Eq.~\eqref{eq:onsager_textbook} \\
Multiplier & $\lambda = 4$ on the right-hand side
  & Sec.~\ref{sec:lambda}; untraced \\
Quantization & Eq.~\eqref{eq:implemented} & --- \\
Orbital moment & $m^z = (e/2\hbar)L^z$, added as $+g\Bfield m^z$
  & Eq.~\eqref{eq:Emod} \\
\bottomrule
\end{tabular}
\caption{Every sign that enters the calculation, and what pins it.}
\end{table}

%======================================================================
\section{Open issues}\label{sec:open}
%======================================================================

\paragraph{1. The multiplier $\lambda = 4$ is untraced.}
This is now the most important caveat in these notes and must be
resolved before publication.  Section~\ref{sec:lambda} establishes
empirically that the right-hand side of Eq.~\eqref{eq:implemented} is
short by a factor of four, and that the factor is flux-independent, but
not \emph{where} the four comes from.  Candidates, none yet checked:
\begin{itemize}
\item The $qq/(2pp)$ primitive-versus-construction-cell flux convention.
  A factor of two applied twice gives four, and the semiclassical
  background field $\Bfield_0 = b_0 (qq/pp)/2$ already carries one such
  factor.  This is the leading suspect.
\item $\phi_0 = h/e$ versus $h/2e$ somewhere in the chain.
\item The shoelace and BZ normalizations that produce $S(E)$ from the
  contour, in \code{isoenergy.py}.
\end{itemize}
Until this is settled, Eq.~\eqref{eq:implemented} should be presented as
a validated form with its validation stated, not as a derivation.

\paragraph{1a. The overall sign in Eq.~\eqref{eq:parity} --- resolved.}
Previous versions of these notes flagged the overall sign as empirical
and in conflict with the traversal argument.  That conflict is closed:
the traversal argument was right and the implementation was wrong.  The
code has been corrected, so Eq.~\eqref{eq:parity} now reads
$\gamma_C = +\sgn(\Bfield)\Phib$, which is simultaneously what the
counterclockwise electron orbit at $\Bfield > 0$ gives by Stokes and what
check (b) selects at all four fluxes.  The three candidate resolutions
formerly listed here --- a mismatch with the textbook form, a
Diophantine effective-field sign, a hole-orbit area convention --- are no
longer needed.  Worth recording, though: the parity signature is only
resolvable at the correct $\lambda$, so the two open questions were
entangled, and a null result on the sign test at $\lambda = 1$ would have
been uninformative rather than negative.

\paragraph{2. Parity of the $dL/dE$ term is unverified.}
As implemented, the orbital-moment correction is even in $\Bfield$
inside the bracket of Eq.~\eqref{eq:implemented}.  If it descends from an
orbit energy shift $\propto m\Bfield$ --- which is odd --- it should
carry a $\sgn(\Bfield)$ exactly as the curvature term does.  This has
\emph{not} been changed, because the term is identically zero in the
non-perturbative channel (Sec.~\ref{sec:bfield}) and the perturbative
channel has no negative-field validation to test against.  Do not trust
perturbative-channel fans at $\Bfield<0$ until this is resolved.

\paragraph{3. Lifshitz threshold is a heuristic.}
See Sec.~\ref{sec:numerics}.  May need retuning in other parameter
regimes.

\paragraph{4. Stale data.}
Fan files generated before the curvature-sign correction and the change
of the $\lambda$ default are wrong in their
\code{SB}/\code{SBM}/\code{SBMC} arrays.  The \code{S}/\code{SM} arrays
are unaffected by the sign but do move with $\lambda$.  Regenerate, or
re-solve from the corresponding \code{*\_detail.mat}.

%======================================================================
\section*{Note on paring down for supplementary information}
%======================================================================

For the SI, the minimal self-contained set is:
Eqs.~\eqref{eq:connection}, \eqref{eq:Oz_kubo}, \eqref{eq:vel_LL},
\eqref{eq:ladder}, \eqref{eq:Phib}, \eqref{eq:parity},
\eqref{eq:implemented}, plus the table in Sec.~\ref{sec:signs} and one
paragraph each on orbit detection and orientation independence
(Sec.~\ref{sec:orientation} is the non-obvious methodological point and
is worth keeping).  Sections~\ref{sec:validation} and \ref{sec:open}
should be compressed to a short validation paragraph --- but issue 1 in
Sec.~\ref{sec:open}, the origin of $\lambda = 4$, needs to be
\emph{resolved}, not merely compressed, before the SI is written.

\end{document}
